% =============================================================
%  ch08_komponenten.tex  --  Komponenten im günstigen Hobby-Bereich
% =============================================================

\chapter{Komponenten im günstigen Hobby-Bereich}

\begin{infobox}[title=Leitprinzip]
Fokus liegt nicht auf Elektrotechnik: Wo möglich Module und Plug-and-Play,
niedrige Kosten. Die Architektur trennt sauber eine
\textbf{Echtzeit-Motorebene} (Mikrocontroller, nah an den Treibern) von einer
\textbf{High-Level-Ebene} (PC/Pi mit vollem ROS\,2 für Trajektorie, Vision,
Kalman).
\end{infobox}

\section{Schrittmotor und Treiber}
\begin{longtable}{@{}p{3.6cm}p{4.2cm}p{7.2cm}@{}}
\toprule
\textbf{Komponente} & \textbf{Empfehlung} & \textbf{Begründung} \\
\midrule
Motor          & NEMA-17 (z.\,B.\ 17HS19) & Standard, günstig, ausreichend Moment mit Getriebe \\
Treiber        & \textbf{TMC2209} (UART)   & StallGuard4 liefert Lastwert über UART = indirekte Drehmomentschätzung; leise, Stepstick-Format, $\approx$5--10\,\euro \\
Treiber-Alt.   & TMC2130/5160 (SPI)       & höherer Strom / SPI statt UART \\
Strom/Spannung & \textbf{INA226}-Modul (I\textsuperscript{2}C) & misst Busspannung + Strom direkt, plug-and-play, $\approx$2\,\euro\ pro Kanal \\
Closed-Loop (opt.) & MKS SERVO42 / BTT S42B & integrierter Encoder + Treiber, falls echter Regelkreis gewünscht \\
\bottomrule
\end{longtable}

\noindent\textbf{Konkret für den Regelungsansatz:} Der TMC2209 im UART-Modus
gibt \code{SG\_RESULT} (Lastproxy) ohne Zusatzhardware. Willst du
\emph{zusätzlich} Busspannung und -strom explizit messen, setzt du je Achse ein
INA226-Modul auf die Treiberversorgung. Beide hängen am selben I\textsuperscript{2}C/UART-Bus
des Mikrocontrollers.

\section{Recheneinheiten}
\begin{description}[style=nextline,leftmargin=1.4em]
\item[ESP32 (Motorebene)] Dual-Core, viele PWM/GPIO, FreeRTOS, harte Echtzeit
  für Schrittgenerierung und das Auslesen von TMC (UART) und INA226
  (I\textsuperscript{2}C). Entscheidend: \textbf{micro-ROS} bindet den ESP32
  als vollwertigen ROS\,2-Knoten in den Graphen ein.
\item[Raspberry Pi (High-Level)] Pi\,4/5 mit Ubuntu kann den vollen
  ROS\,2-Stack tragen (Trajektorienplanung, Kalman). Der \textbf{Pi Zero} ist
  dafür zu schwach -- für Vision ungeeignet; nutze ihn höchstens als
  micro-ROS-Brücke. Vision und schweres Rechnen gehören auf den externen PC.
\end{description}

\begin{infobox}[title=Empfohlene Arbeitsteilung]
\textbf{Externer PC} (Ubuntu + ROS\,2): Trajektorienplanung, Stereo-Vision,
Kalman-Filter, Visualisierung.
\;$\longleftrightarrow$\; (WLAN/USB) \;$\longleftrightarrow$\;
\textbf{ESP32} (micro-ROS): empfängt Sollwerte, generiert Schritte, liest
StallGuard/INA226, publiziert Gelenkzustände. Die Stereokamera hängt per USB
am PC.
\end{infobox}

\section{Stereokamera}
Das ELP-Modell ist eine synchronisierte Doppelobjektiv-USB-Kamera
(UVC-konform). Sie meldet sich unter Linux als V4L2-Gerät (\code{/dev/video*})
und liefert in der Regel ein \emph{nebeneinanderliegendes} (side-by-side) Bild
über eine USB-Verbindung -- die Synchronisierung beider Sensoren ist der
Vorteil gegenüber zwei Einzelkameras (wichtig für bewegte Objekte wie den Ball).
In ROS\,2:
\begin{itemize}
\item Treiber: \code{v4l2\_camera} oder \code{usb\_cam}.
\item Side-by-side-Frame in linkes/rechtes Bild splitten (kleiner eigener Knoten,
  siehe Teil~VI).
\item Eigene Stereokalibrierung ist Pflicht -- das Modul ist nicht vorkalibriert.
\end{itemize}
